\documentclass[11pt,a4paper]{article}
\usepackage[utf8]{inputenc}
\usepackage[T1]{fontenc}
\usepackage[spanish]{babel}
\usepackage{hyperref}
\usepackage{geometry}
\geometry{margin=2.5cm}
\usepackage{parskip}

\title{348 verbalizar tareas y verbalizar procesos no son lo mismo}
\author{Alejandro Valenzuela Rivera\\\small Universidad Autónoma de Nuevo León}
\date{\today}
\begin{document}
\maketitle

\begin{quote}
\textbf{Sesgo declarado.} Este trabajo se sostiene desde el sesgo declarado de su autor, y no finge neutralidad. Lo que afirma se puede tumbar, y ese es el punto: el falsificador está escrito, no escondido.
\end{quote}

\section*{Resumen}
Una cosa es verbalizar un proceso y otra verbalizar una tarea: el cuerpo ejecuta tareas, la mente las abstrae, y los procesos son lo que no sabemos.

\section{Por qué se sostiene}
Si la tarea y el proceso se verbalizan igual, el ente no distingue lo que hace de lo que no entiende.

\section{De dónde viene}
La autoridad: 'el cuerpo ejecuta tareas, la mente las abstrae y los procesos son lo que no sabemos'.

\section{El caso}
El ente: las tareas viven en el cuerpo, y los procesos son el hueco que nadie escribió.

\section{Método}
\textquestiondown{}La marca distingue lo que se ejecuta de lo que no se sabe?

\section{Límites, y qué lo tumba}
Un ente donde la tarea y el proceso lleven la misma marca.

\section{Lo que queda abierto}
\textquestiondown{}Se verbalizó la tarea, o el proceso que la hace aparecer?

\section*{Cómo se leyó esta tesis}
Esto se escribió en el repositorio man-made-empathy, del que este candado es una pieza: candados/01 el ente y el metodo/348 verbalizar tareas y verbalizar procesos no son lo mismo.md. La máquina que lo compuso lo hizo desde este mismo archivo, y no agregó nada que no estuviera escrito. Si algo de esto es falso, se dice en el mismo sitio, y se corrige con fecha.

\end{document}
